\section*{Chapitre \ref{ch:hf:bond-and-etf-request-for-quote-market-making} --- Market making sur demandes de cotation (RFQ) d'obligations et d'ETF}

\begin{solution}{exo:hf:bond-and-etf-request-for-quote-market-making:1}
$100\times(1-0.0006)-1\,250/50\,000=99.94-0.025=\$99.915$.
\end{solution}

\begin{solution}{exo:hf:bond-and-etf-request-for-quote-market-making:2}
$25\times1.163=29.1$ cents.
\end{solution}

\begin{solution}{exo:hf:bond-and-etf-request-for-quote-market-making:3}
$1/(1+e^{-(0.665-0.0612\times42.5)})=12.6~\%$~; à 20 cents, $1/(1+e^{0.559})=36.4~\%$.
\end{solution}

\begin{solution}{exo:hf:bond-and-etf-request-for-quote-market-making:4}
Avec davantage de dealers, l'estimation du vainqueur est plus sélectionnée~: son erreur espérée sachant qu'il gagne augmente (de 35 à 44 cents)~; le markup doit la couvrir, et
au-delà de quelques dealers, cet effet l'emporte sur la concurrence.
\end{solution}

\begin{solution}{exo:hf:bond-and-etf-request-for-quote-market-making:5}
Le dealer n'apprend la meilleure cotation des concurrents que lorsqu'il gagne, c'est-à-dire lorsque ceux-ci cotaient bas~; les pertes indiquent seulement que quelqu'un cotait plus haut.
Un modèle calibré sur les seuls \omterm{def:hf:bond-and-etf-request-for-quote-market-making:cover}{prix de couverture} sous-estime la concurrence~; les pertes doivent entrer comme observations censurées.
\end{solution}

\begin{solution}{exo:hf:bond-and-etf-request-for-quote-market-making:6}
Les erreurs du dealer sur des centaines de lignes se compensent en partie~; la valeur du portefeuille est estimée plus précisément que chaque ligne, si bien que l'effet de sélection sur
l'ensemble est plus faible.
\end{solution}

\begin{solution}{exo:hf:bond-and-etf-request-for-quote-market-making:7}
Avec une erreur de 15 cents face aux 25 cents des concurrents, notre meilleur markup est de 25 cents~; nous gagnons 2.48 cents par demande et remportons 23.4~\% des demandes~: une meilleure
estimation rapporte à la fois du volume et de la marge.
\end{solution}

\begin{solution}{exo:hf:bond-and-etf-request-for-quote-market-making:8}
Le profit sur une victoire est le markup diminué de l'erreur d'estimation, et l'erreur conditionnelle à une victoire est positive et importante (35 cents avec cinq dealers dans le
modèle)~: à un markup de 25 cents, le cotateur perd en moyenne.
\end{solution}

\begin{solution}{pb:hf:bond-and-etf-request-for-quote-market-making:1}
\begin{enumerate}
\item Ses dernières transactions, devenues anciennes, sont ramenées à un composite déplacé par la courbe de l'émetteur, des comparables avec leurs sensibilités et des proxies liquides.
\item Voir la \cref{def:hf:bond-and-etf-request-for-quote-market-making:auto}.
\item Les enchères électroniques unilatérales sont viables même pour des obligations peu actives~; la négociation électronique par RFQ, encore fragmentée, modifie les coûts et la
qualité d'exécution sur l'ensemble du marché.
\item MarketAxess en 2025~: chiffre d'affaires record de 846~millions de dollars, transactions de portefeuille en hausse de 48~\% à 1.4~milliard de dollars par jour, blocs en hausse de 24~\% à 5~milliards de dollars.
\item Voir la \cref{def:hf:bond-and-etf-request-for-quote-market-making:cover} et la \cref{def:hf:bond-and-etf-request-for-quote-market-making:win}.
\item $P(\text{victoire}\mid m)\,\mathrm{E}[\text{bénéfice}\mid\text{victoire},m]$~; le second terme est diminué de l'erreur de l'estimation qui a gagné, laquelle est
positive.
\item 29 et 38.5 cents.
\item Meilleures réponses à un markup commun des concurrents, jusqu'à ce que le markup se reproduise lui-même.
\item 47.5 cents, 7.96 cents, 27.7~\%~; 42.5, 0.94, 13.2~\%~; 50.0, 0.45, 6.9~\%.
\item Il cote 20 à 25 cents plus serré, gagne plus de deux fois plus souvent, et perd à partir de cinq dealers.
\item 12.6~\% ajusté contre 13.2~\% simulé au markup d'équilibre.
\item 99.915~\$ par action.
\item Le profit par demande passe de 7.96 cents seul à 0.94 avec cinq dealers et à environ 0.4 avec huit à dix~; le taux de victoire passe de 28~\% à 13~\% et
à 7~\%~; le markup reste entre 40 et 50 cents.
\item Assez pour susciter la concurrence, peu nombreux pour limiter la malédiction du vainqueur et l'information divulguée (\texttt{firm.rfq.best\_n} du livre~2).
\item Taux de victoire et marge, les deux~: 23.4~\% et 2.48 cents par demande face à cinq concurrents moyens, contre 13.2~\% et 0.94.
\item La diversification réduit la malédiction sur un portefeuille, et les fonds offrent une couverture.
\item Les demandeurs anonymes ne peuvent pas être hiérarchisés~; le système de \omterm{def:hf:bond-and-etf-request-for-quote-market-making:auto}{cotation automatique} est en concurrence avec les investisseurs.
\item Le répondeur sur obligations d'entreprise.
\item Comme observations censurées~: les \omterm{def:hf:bond-and-etf-request-for-quote-market-making:cover}{prix de couverture} bornent les concurrents par le bas lorsqu'on gagne, les pertes par le haut.
\item Un prix, en moins d'une seconde, pour une obligation que le marché n'a pas valorisée depuis des semaines.
\end{enumerate}
\end{solution}

\begin{solution}{iq:hf:bond-and-etf-request-for-quote-market-making:1}
Vous avez payé 30 cents de plus que nécessaire pour gagner~; soit votre estimation était élevée, soit votre markup trop mince. Une observation isolée est du bruit~; de nombreuses alimentent le
\omterm{def:hf:bond-and-etf-request-for-quote-market-making:win}{modèle de probabilité de gain}.
\iqlookfor{le \omterm{def:hf:bond-and-etf-request-for-quote-market-making:cover}{prix de couverture} comme information, et la malédiction du vainqueur.}
\end{solution}

\begin{solution}{iq:hf:bond-and-etf-request-for-quote-market-making:2}
Dernières transactions déplacées par la courbe de l'émetteur depuis, spreads des comparables ajustés de l'échéance et du rang de séniorité, et mouvement du marché tiré des proxies liquides~; élargir
l'incertitude avec le temps écoulé depuis la dernière transaction.
\iqlookfor{les comparables et l'ancienneté de l'information.}
\end{solution}

\begin{solution}{iq:hf:bond-and-etf-request-for-quote-market-making:3}
Des fair values précalculées et mises à jour sur les mouvements de marché~; un \omterm{def:hf:bond-and-etf-request-for-quote-market-making:win}{modèle de probabilité de gain} rapide~; des limites sur la taille, l'inventaire et les listes d'obligations~; l'acheminement des demandes inhabituelles vers les traders.
Risques~: intrants obsolètes, données de référence erronées, cotation incontrôlée en situation de stress.
\iqlookfor{précalcul et limites.}
\end{solution}

\begin{solution}{iq:hf:bond-and-etf-request-for-quote-market-making:4}
Vérifier si les fair values ont dérivé à la hausse (la malédiction), si des concurrents se sont retirés, si la composition de la clientèle a changé~; le P\&L des gains récents face aux
prix ultérieurs.
\iqlookfor{le taux de victoire comme signal d'alerte.}
\end{solution}

\begin{solution}{iq:hf:bond-and-etf-request-for-quote-market-making:5}
Avec vingt dealers, la malédiction est forte~: il faut coter plus large, ou ne pas coter du tout, sauf si l'estimation est exceptionnellement bonne ou si l'inventaire appelle l'obligation.
\iqlookfor{la malédiction croît avec le nombre de dealers interrogés.}
\end{solution}

\begin{solution}{iq:hf:bond-and-etf-request-for-quote-market-making:6}
Le maximum de deux variables $N(0,\sigma^2)$ indépendantes a pour espérance $\sigma/\sqrt\pi\approx0.564\sigma$.
\iqlookfor{statistiques d'ordre.}
\end{solution}
